\section{Time and the Arrow}

We have space. Now time. We touched it in the chapter on beats. Here we make it physical and answer the question physics finds hardest: why does time have a direction, when the underlying laws do not seem to care which way it runs?

\subsection{Time is the depth of cause}

Space was the sideways shape of the web, how many notes separate two vibes. Time is the other direction through the web, the depth of before-and-after. A vibe's state was set from its neighbors' earlier states, which were set from theirs, and so on back. Those chains of ``this came from that'' are the threads of time. How much time has passed to reach some event is, in effect, how long the longest such chain leading up to it is.

So time and space are the same web read two ways: across it (who is beside whom) is space, and along the chains of causation (what came from what) is time. This is exactly the deep lesson of relativity, that space and time are facets of one fabric, and here it is built in from the start because both are just structure in the one web.

\subsection{Why time has an arrow}

Here is the famous puzzle. The basic laws of physics look the same run forwards or backwards. A film of two billiard balls colliding looks fine played in reverse. Yet the world has an unmistakable arrow: eggs break but never unbreak, heat flows from hot to cold and never back, we remember the past and not the future. Where does that one-way arrow come from, if the laws are two-way?

In this theory the arrow is not added on. It is the growth. The crystal only ever gets bigger. New vibes are added at the frontier, never removed. Notes accumulate, never undone. The heap of what-has-happened only grows. That is the arrow, plain and irreversible, and everything else lines up behind it: you remember the past because the past is the fixed interior recorded in the structure, and you cannot remember the future because it has not been made yet.

We saw the root of this in the bootstrap. Distinctions can be added but never subtracted, because erasing one would itself be a new distinction. There is no road back. So the direction of time is free, falling straight out of the fact that the universe is a growing structure rather than a static block.

\subsection{The present is the edge, again}

This gives a picture of ``now'' quite different from the everyday one. The present is not a thin slice sweeping through a finished block of spacetime. The future block does not exist yet to be swept through. The present is the living, growing rim of the crystal, the freshly added cells where the action is. The past is everything already grown, fixed behind. The future is the open room, always waiting in hyperbolic space, for the next cells.

You are not a traveler moving along a pre-laid track of time. You are part of the edge where time is being made, beat by beat. That is why the future feels open and the past feels fixed: they are, literally, the unmade and the made parts of one growing crystal.

\subsection{No master clock, no paradox}

Recall there is no universal clock, because a single cosmic ``now'' would be a forbidden preferred frame. Different observers slice the before-and-after web differently, and none is privileged, exactly as relativity demands. Yet the arrow survives this, because the arrow is not about any particular slicing. It is about growth versus shrinkage, and the crystal only grows, for everyone. So time can be relative in its slicing and absolute in its direction, with no contradiction.

\textbf{Takeaway:} time is the depth of the before-and-after web, the same fabric as space read along the chains of cause. Its arrow is not added but is the one-way growth of the crystal, distinctions only piling up. The present is the growing edge, the future the open room ahead, and the arrow survives the absence of any universal clock.
